% =====================================================================
\appendix

\section{Glossary}
\label{app:glossary}
{\small
\begin{longtable}{@{}L{3.4cm} L{7.1cm} L{5.2cm}@{}}
\toprule Term & Plain-language meaning & Why it matters here\\ \midrule\endhead
\bottomrule\endfoot
Quantum advantage & Better cost or capability for a stated output, accuracy and classical comparator & There is no unqualified winner\\
Physical and logical qubits & Many physical devices are combined to protect one logical qubit against errors & Resource counts differ by orders of magnitude\\
Fault tolerance & Error correction that sustains a long computation at a stated failure rate & Distinct from today's error-mitigated experiments\\
Error mitigation & Estimating a less noisy result from extra noisy runs and a noise model & Bias and measurement overhead matter\\
Phase estimation (QPE) & A quantum procedure that estimates energies from a prepared state's spectral content & One run need not find the wanted state\\
Spectral weight, overlap & How much of the prepared state lies in the wanted energy region & Sets the probability of a useful run\\
One-shot and expected runtime & Time for one execution; time including repeats after detected failure & Retries must be counted once\\
Magic-state factory & An error-corrected subroutine that supplies the resource for costly logical gates & Can occupy most physical qubits\\
T, Toffoli, CCZ & Different costly gate types used in compilation & Counts are not interchangeable\\
Code cycle, reaction time & The clock of error correction; the delay of the classical controller & Set how fast logical gates can run\\
Tensor network & A compressed representation of a many-body state or computation that exploits structure & Classical cost depends on geometry and entanglement\\
DMRG, bond dimension & A variational tensor-network solver; the bond dimension sets the size of the representation & Accuracy and cost depend on convergence in it\\
Neural quantum state & A trainable wavefunction optimized directly for a Hamiltonian & A direct solver, not a surrogate trained on labels\\
Sign problem, sign-free & Statistical cancellation that makes Monte Carlo costly; absent in special cases & Changes the classical baseline\\
Active space & The electrons and orbitals treated explicitly in a chemistry model & A model result is not the molecule or enzyme\\
Chemical accuracy & An energy scale near 1\,kcal/mol, about 1.6\,mHa & The observable may need another tolerance\\
XEB & A score comparing measured outputs with ideal circuit probabilities & Matching it is not simulating the distribution\\
OTOC, Loschmidt echo & Diagnostics of how quantum dynamics spreads or reverses information & The observable, not the full state, is the task\\
Amortization & Paying a data and training cost once and reusing the result & Workload size changes the comparison\\
Condition number & Ratio of a matrix's largest to smallest singular value & Sets the cost of quantum linear solvers; grows with mesh refinement\\
HHL, QSVT & Quantum algorithms that prepare a state proportional to $A^{-1}b$ & Fast in size, slow in condition number and readout\\
Block encoding & Embedding a matrix inside a quantum circuit so that it can be applied & Its cost and scaling factor enter the total\\
Carleman linearization & Rewriting a nonlinear equation as a larger linear one & Efficient only for weak nonlinearity\\
Multigrid & A classical solver whose cost is proportional to the number of unknowns & The incumbent for elliptic problems\\
Query complexity & Number of input lookups, with other computation free & Does not determine total runtime\\
\end{longtable}
}

% =====================================================================
\section{Evidence entries}
\label{app:entries}
Table~\ref{tab:master} lists every entry. An entry is a pairing: the same Hamiltonian, circuit family, molecule, material or optimization class addressed on quantum hardware or in a fault-tolerant estimate, and by a classical method; or a theoretical result that bears directly on such a pair; or, as context, an AI result commonly cited as making quantum computing unnecessary. Stated assumptions: agreement is categorized by the kind of match (Section~\ref{sec:match}); fault-tolerant estimates are not yet subject to engineering losses; roadmap figures are announced targets; and where only parallel progress exists we infer no winner.

{\scriptsize
\setlength{\tabcolsep}{3pt}
\begin{longtable}{@{}r L{0.7cm} L{2.25cm} L{1.9cm} L{3.1cm} L{3.4cm} L{1.85cm} L{1.3cm}@{}}
\caption{The \nEntries{} evidence entries (status as of 11 November 2026). Type: E1 device comparison, E2 resource estimate, E3 theory, E4 context (Section~\ref{sec:types}). Method class: AI = learned; CV = classical variational; SC = structured numerical. ``Resource'' gives the coverage of the linked problem class in Table~\ref{tab:forward}; ``---'' where the entry has no forward task. Entry numbers are stable identifiers (\#$n$ is E$n$ in the ledger). Dates of claims and responses, tolerances and workloads are in \texttt{entries.csv}.}\label{tab:master}\\
\toprule
\# & Type & Problem & Quantum side & Classical/AI counter (class) & Accuracy and cost & Present result & Resource\\
\midrule\endfirsthead
\toprule
\# & Type & Problem & Quantum side & Classical/AI counter (class) & Accuracy and cost & Present result & Resource\\
\midrule\endhead
\bottomrule\endfoot
\multicolumn{8}{@{}l}{\textit{Quantum many-body physics and optics}}\\
1 & E1 & Random circuit sampling, 53\,q & Google Sycamore, 2019 & Tensor-network contraction (SC)~\cite{pan2022,hangleiter2026} & Samples at comparable XEB; later RCS experiments not entered individually & Validation benchmark matched & ---\\
2 & E1 & Gaussian boson sampling, Jiuzhang 2.0 & USTC, 2021~\cite{zhong2021} & Marginal-based mock-up samplers~\cite{villalonga2021}; later a tensor-network algorithm~\cite{oh2024} (SC) & Mock-up closer than the experiment to the ideal marginals (\eld{E02} d); benchmark statistics outperformed in 2023; quantum model favoured over a squashed-state alternative~\cite{martinez2023} & Validation statistic challenged, later matched & ---\\
3 & E1 & Kicked Ising, 127\,q, 20 steps & IBM Eagle, 2023 & BP tensor networks~\cite{tindall2024}; sparse Pauli dynamics~\cite{begusic2023,begusic2024} (SC) & Ideal-circuit observables within or below experimental error; one laptop core & Observables reproduced & class \#4: \mtP\\
4 & E1 & Spin-glass annealing dynamics, 2D/3D & D-Wave Advantage2~\cite{king2025} & BP-TN~\cite{tindall2025} (SC); Jastrow--Feenberg t-VMC to 128 spins~\cite{mauron2025} (CV) & Lower error than QPU in some regimes; not all observables~\cite{kingreply} & Partly reproduced; \psC & \mtN\\
5 & E1 & 2D TFIM Floquet prethermalization, 56\,q & Quantinuum H2~\cite{haghshenas2025} & TN, NN, operator truncation (SC+AI) & No tested method converged in the hardest regime & \psQ{} (\eld{E05} d) & \mtN\\
6 & E1 & OTOC(2), 65\,q & Google Willow~\cite{google2025otoc} & TN contraction, TNMC, TNBP (SC)~\cite{bermejo2026} & 2.1\,h of data collection per circuit (measured) against 3.2\,y on Frontier per circuit (cost model); accuracy at 65 qubits from an error model & \psQ{} (\eld{E06} d) & \mtN\\
7 & E1 & Operator Loschmidt echo, 56\,q & IBM Heron + Algorithmiq~\cite{barron2026} & Three classical groups (SC) & Classical predictions mutually inconsistent & \psQ{} (preprint) & \mtN\\
8 & E4 & Rydberg spin-liquid diagnostics & Harvard simulator, 2021 & Time-dependent VMC (CV)~\cite{mauron2025tee} & R\'enyi TEE $\approx$0.48 vs $\ln2$; topological order without RVB TEE & Comple\-mentary inference & \mtN\\
\multicolumn{8}{@{}l}{\textit{Quantum chemistry and electronic structure}}\\
9 & E1 & Fe--S clusters by SQD & IBM Heron + Fugaku~\cite{shirakawa2025} & Selected CI; single-layer UCJ energy evaluation~\cite{belagali2026}; GPU DMRG~\cite{legeza2026} (SC); ML selectors (AI) & Ideal QSCI 50{,}551 vs HCI 4{,}260 dets at 0.1\,Ha error~\cite{reinholdt2025}; single-layer [4Fe-4S] energy on a laptop $<$1\,min; CAS(54,36) DMRG in \feDmrgH\,h on one GPU node & No advantage claimed; \psM{} for the single-layer energy objective & ---\\
10 & E2 & FeMoco 76-orbital model & THC~\cite{lee2021thc}; dense code~\cite{low2026dense} & UDMRG + UCCSDTQ~\cite{zhai2026} (SC) & $\pm$0.31 (consensus) / $\pm$1.25 (conservative) kcal/mol, extrapolation-dependent; quantum \femoT\,h expected per accepted run & Classical feasibility; operational ranking unresolved & class \#10: \mtP\\
11 & E3 & Generic ground-state chemistry & QPE hypothesis & Empirical complexity of classical heuristics~\cite{lee2023} (SC) & No evidence of generic exponential separation in studied classes & Theory & \mtN\\
12 & E4 & Main-group thermochemistry & --- & Skala XC functional~\cite{luise2025} (AI) & W4-17 at chemical accuracy; GMTKN55 WTMAD-2 3.89 (v1.0), 2.72 (v1.1) kcal/mol & No quantum comparator & ---\\
\multicolumn{8}{@{}l}{\textit{Materials}}\\
13 & E4 & Solid-electrolyte discovery & --- & ML+DFT+MD on $\sim$1{,}000 cloud VMs (AI)~\cite{chen2024} & 32.6\,M $\to$ 18 candidates in $<$80\,h; Na$_x$Li$_{3-x}$YCl$_6$ synthesized, conductivity measured & Validated classical result & ---\\
14 & E4 & Crystal stability at scale & --- & GNoME~\cite{merchant2023} (AI) & 2.2\,M stable vs prior data; $\sim$381{,}000 on updated hull & No quantum comparator & ---\\
\multicolumn{8}{@{}l}{\textit{Life sciences}}\\
15 & E1 & Lattice protein folding & VQE, 9--127\,q~\cite{robert2021,doga2024,zhang2026} & Exhaustive/Gurobi on same Hamiltonian (SC) & RMSD 1.781\,\AA{} (quantum best) vs 1.879\,\AA{} (classical model optimum) vs 3.53\,\AA{} (AF2) & Objectives differ; no computational advantage shown & ---\\
16 & E1 & KRAS inhibitor generation & QCBM+LSTM, IBM~\cite{kras2025} & Classical generative priors (AI) & Modest benchmark gain; no superiority claimed~\cite{uoft2025} & Quantum contribution unisolated & \mtN\\
\multicolumn{8}{@{}l}{\textit{Optimization}}\\
17 & E4 & Spin-glass ground states & Simulated QA (emulation) & Variational neural annealing~\cite{hibat2021} (AI) & Beats SA and SQA over studied sizes & No hardware comparator & ---\\
18 & E1 & Approx.\ optimization, 2D spin glass & D-Wave + QAC~\cite{munoz2025} & PT-ICM; simulated bifurcation~\cite{pawlowski2026} (SC) & QAC scaling advantage vs PT-ICM at gap $\ge$1\%; follow-up finds no clear advantage vs SBM & \psC & \mtN\\
27 & E1 & Maximum independent set, $\le$289\,q & Harvard Rydberg array~\cite{ebadi2022} & Exact and heuristic solvers; relaxed SA (SC)~\cite{andrist2023} & Superlinear speedup claimed vs SA on hardest graphs; quasi-planar instances solved to optimality at thousands of nodes in minutes & \psC & \mtN\\
19 & E4 & MIS on $10^6$-node graphs (control) & --- & Physics-inspired GNN vs greedy~\cite{angelini2022,schuetz2023} & Greedy $10^4\times$ faster and better on sparse MIS & Benchmarking control & ---\\
\multicolumn{8}{@{}l}{\textit{Machine learning theory}}\\
20 & E3 & Recommendation, PCA, low-rank SDP & Quantum linear algebra & Dequantized sampling~\cite{tang2023,tang2021} (SC) & Polynomial slowdown under $\ell^2$ sampling access & Theory & \mtN\\
21 & E3 & Ground-state properties within a phase & QC simulation target & Classical ML on classical shadows~\cite{huang2022} (AI) & Efficient after data access, gapped within-phase assumptions & Theory & ---\\
\multicolumn{8}{@{}l}{\textit{Nuclear physics}}\\
22 & E4 & Nuclear ground states, $A\le20$ & Few-nucleon QC demos & NQS VMC to $^{16}$O and $A{=}20$~\cite{lovato2022,gnech2024,lovato2026} (AI) & Energies, radii, magnetic moments for cited interactions & No matched comparator & ---\\
\multicolumn{8}{@{}l}{\textit{Further device, workflow and theory entries}}\\
23 & E1 & Heavy-hex Floquet magnetization, $\le$74\,q & IBM Heron + QESEM~\cite{leviatan2026} & Tensor networks; sparse Pauli-path on GPUs and Fugaku (SC) & Authors: TN not converged, Pauli-path truncation-dependent; unbiased and extrapolated estimators agree where both exist; cross-platform checks at selected cycles & \psQ{} (\eld{E23} d) & \mtN\\
24 & E1 & Doped Clifford sampling, 70\,q (original) & IBM Heron, spacetime codes~\cite{martiel2026} & Tensor-network contraction (SC)~\cite{manabe2026} & Amplitudes of the 2{,}051 published bitstring batches in \dcMin\,min on 256 H100 GPUs & Validation benchmark matched (16 d) & ---\\
25 & E1 & Doped Clifford sampling, 64\,q (revised) & IBM Heron; depth 73, 314 T, periodic boundaries~\cite{martiel2026} & None found & Fidelity lower bound 0.349 (95\%) per authors; periodic boundaries remove the contraction path of \#24 & \psQ{} (\eld{E25} d) & ---\\
26 & E1 & Driven Bose--Hubbard sampling, 64 sites, 20 atoms (version of April 2026) & USTC quantum gas microscope~\cite{zheng2026} & Tensor networks (authors' comparison) & Authors: sampling rate $10^3\times$ a supercomputer; validated by Bayesian tests & \psQ{} (\eld{E26} d) & ---\\
28 & E1 & Embedded protein--ligand electronic structure & Two IBM Heron r2, $\le$94\,q; \merzQpu\,QPU-h~\cite{merz2026} & SCI, DMRG, CCSD on fragments (SC) & Fragment errors 0.10--2.3\,mHa vs DMRG; SCI better on one fragment; no sampler ablation & No advantage claimed & ---\\
29 & E4 & Frustrated and doped lattice models & --- & Transformer NQS~\cite{rende2024,gu2026} (AI) & $J_1$--$J_2$ $10\times10$: $E/N=-0.497634(1)$, below DMRG; 4 days on 20 A100 GPUs & No matched comparator & classes \#6, 36\\
30 & E4 & Molecular excited states & --- & Neural-network VMC~\cite{pfau2024,schatzle2026} (AI) & Excited states and transition observables; numbers not extracted & No matched comparator & ---\\
31 & E3 & Unstructured decision problems, query model & Any $T$-query algorithm & Classical queries to influential variables~\cite{aa2014,turturean2026} (SC) & $\le(4T)^8\epsilon^{-14}\delta^{-8}$ classical queries on all but a $\delta$ fraction of inputs & Theory; claimed proof, machine-checked core & class \#29\\
32 & E3 & Engineering PDEs by quantum linear and differential-equation solvers & HHL/QSVT; quantum ODE solvers; Carleman linearization & Multigrid and Krylov solvers; explicit time stepping; Monte Carlo (SC)~\cite{montanaro2016,linden2022,an2025,lewis2024} & At most polynomial speedup in fixed dimension; overheads from dissipation, non-normality and strong nonlinearity & Theory & class \#38\\
33 & E2 & Drag force by quantum lattice-Boltzmann solvers & Resource estimates~\cite{penuel2025,jennings2025,zhuang2025} & Explicit classical CFD, $O(\mathrm{Re}^3)$ (SC) & $10^{21}$--$10^{39}$ logical qubits $\times$ T gates for a sphere; bounded gain $\mathrm{Re}^{3D/8}$; one crossover reported at $2^{80}$ grid points (source-reported range; public records reach $7.9\times10^{40}$) & No practical advantage established; conditional polynomial proposals & class \#38: \mtP\\
34 & E1 & Gaussian boson sampling, Jiuzhang (2020), 100 modes & USTC, 2020~\cite{zhong2020} & Marginal-based mock-up samplers~\cite{villalonga2021}; squashed-state model~\cite{martinez2023} (SC) & Mock-up closer than the experiment to the ideal marginals (\eld{E34} d); a classical explanation not ruled out; no sampling of the full ideal distribution & Validation statistic challenged & ---\\
35 & E1 & Gaussian boson sampling, Borealis & Xanadu, 2022~\cite{madsen2022} & Tensor-network GBS algorithm~\cite{oh2024} (SC) & Benchmark statistics of the experiment outperformed (\eld{E35} d) & Validation benchmark matched & ---\\
36 & E1 & Gaussian boson sampling, Jiuzhang 3.0 & USTC, 2023~\cite{deng2023} & Tensor-network GBS algorithm~\cite{oh2024} (SC) & Benchmark statistics of the experiment outperformed (\eld{E36} d); newest generation open~\cite{hangleiter2026} & Validation benchmark matched & ---\\
\end{longtable}
}


% =====================================================================
\section{Forward records}
\label{app:forward}
Each of the \nClasses{} problem classes has two independent records, present comparative status and resource coverage (Section~\ref{sec:types}), and three statements of the evidence that would change it. A class takes the least mature coverage label among the records its statement depends on. ``Not applicable'' means that no resource-comparison label is assigned to that record; it does not mean that resource accounting is unnecessary. A trigger has a fixed target: moving from half filling to doping, from a proteome to a fragment, or from a full field to one scalar opens a new record and does not reverse the old one. A classical estimate that fails reopens a comparison; it does not establish a quantum result. The six screens of Table~\ref{tab:checklist} organize the evidence; they are not sequential gates and none is an exclusion proof. The classes are not independent units: several share sources, Hamiltonians or nested scopes, linked by the task-family field of the ledger, and classes 34 and 35 are placed in Fig.~\ref{fig:grid} by their full-field sub-task.

\begin{table}[htbp]
\centering\scriptsize
\setlength{\tabcolsep}{4pt}
\caption{Screening checklist. ``Unknown'' is always a permitted answer. Answers inform the two records of Fig.~\ref{fig:grid}; they do not by themselves decide a class.}\label{tab:checklist}
\begin{tabular}{@{}L{3.1cm} >{\columncolor{qcF!10}}L{3.0cm} >{\columncolor{qcNl!40}}L{3.0cm} >{\columncolor{qcUl!60}}L{2.5cm} L{3.2cm}@{}}
\toprule
Screen & Supports quantum opportunity & Supports classical capability & Unknown & Evidence needed\\\midrule
S1 Speedup on the actual instance family & Exponential or structured super-quadratic & Unstructured decision problem (polynomial query limit; quadratic search under the Sec.~\ref{sec:breakeven} model) & Speedup unproven for these instances & Algorithm analysis on representative instances\\
S2 Input and output & Few scalars or observables out; native quantum input & Bulk classical data in; full-state readout from an amplitude-encoded solution & Access model unstated & Loading model; estimator and shot count\\
S3 Workload & One or few hard instances & Many moderate queries a validated surrogate amortizes & $K$ and surrogate adequacy unknown & $T_C(K)$, $T_Q(K)$ (Eqs.~\ref{eq:tck}--\ref{eq:tqk})\\
S4 Adequacy & The question still needs this computation & A validated reformulation answers it & Validation domain unclear & Error budget (Eq.~\ref{eq:errbudget})\\
S5 Classical structure & No tested classical method converges at relevant sizes & A structured, variational or neural method reaches the tolerance (incl.\ sign-free sectors) & Classical scaling not established & Converged SC, CV and AI baselines on the same instance\\
S6 Resource evidence & Coverage \mtC{} with validated assumptions & --- & Partial, scenario or missing; assumptions unvalidated & Execution record (Appendix~\ref{app:models})\\
\bottomrule
\end{tabular}
\end{table}

\begin{figure}[htbp]
\centering
\resizebox{\linewidth}{!}{\begin{tikzpicture}[gridcell/.style={font=\footnotesize\bfseries, text width=2.0cm, align=center, inner sep=1pt}]
% background cells
\fill[qcF!20] (0,0) rectangle (15.4,-1.05);
\fill[qcF!14] (0,-1.05) rectangle (15.4,-2.1);
\fill[qcF!7] (0,-2.1) rectangle (15.4,-3.15);
\foreach \i in {0,...,6} {\foreach \j in {0,...,4} {
  \draw[black!25] (\i*2.2,-\j*1.05) rectangle ++(2.2,-1.05);
}}
% generated by ledger/make_figdata.py
\node[gridcell] at (1.100,-3.675) {1, 2, 3};
\node[gridcell] at (7.700,-1.575) {4, 6, 7, 10, 11, 12, 20};
\node[gridcell] at (5.500,-3.675) {5, 30};
\node[gridcell] at (7.700,-3.675) {8, 13, 14, 21, 23, 28, 31, 36, 37};
\node[gridcell] at (7.700,-0.525) {9};
\node[gridcell] at (3.300,-4.725) {15, 26};
\node[gridcell] at (7.700,-2.625) {16};
\node[gridcell] at (14.300,-4.725) {17};
\node[gridcell] at (12.100,-4.725) {18};
\node[gridcell] at (9.900,-4.725) {19, 22, 24, 25, 27, 32};
\node[gridcell] at (3.300,-2.625) {29};
\node[gridcell] at (9.900,-3.675) {33, 35};
\node[gridcell] at (9.900,-2.625) {34};
\node[gridcell] at (3.300,-1.575) {38};

% axis labels
\foreach \lab [count=\i from 0] in {standing claim, theoretical result, contested, not compared, validated classical capability, matched classical result (scoped), capability benchmark} {
  \node[font=\scriptsize, text width=2.05cm, align=center, anchor=north] at (\i*2.2+1.1,-5.32) {\lab};
}
\foreach \lab [count=\j from 0] in {complete, partial, scenario, none, not applicable} {
  \node[font=\scriptsize, anchor=east] at (-0.1,-\j*1.05-0.525) {\lab};
}
\node[font=\footnotesize\bfseries, anchor=north] at (7.7,-6.25) {Present comparative status};
\node[font=\footnotesize\bfseries, rotate=90, anchor=south] at (-1.75,-2.625) {Resource coverage};
\end{tikzpicture}}
\caption[Forward classes on two independent axes]{The \nClasses{} problem classes (numbers as in Table~\ref{tab:forward}) placed on two independent axes. The shading only separates the three coverage rows and carries no other meaning. Cell contents are generated from \texttt{forward\_classes.csv}. The ``not compared / partial'' cell holds resource-estimated candidates whose classical comparison is unresolved. The single entry in the top row is complete in coverage at the logical level, with unvalidated assumptions. No column means ``classical lead'': ``validated classical capability'' records a demonstrated classical result with no matched quantum workflow. Classes 34--35 are placed by their full-field sub-task. The classes are not independent units; several share sources, Hamiltonians or nested scopes, linked by the task-family field of the ledger, so the cell counts are counts of coded labels and not a survey statistic.}
\label{fig:grid}
\end{figure}

{\scriptsize
\setlength{\tabcolsep}{3pt}
\begin{longtable}{@{}r L{2.4cm} L{2.4cm} L{3.1cm} L{2.3cm} L{2.4cm} L{1.9cm}@{}}
\caption{Forward records for \nClasses{} problem classes (\texttt{forward\_classes.csv}). ``Resource evidence'' is in original units with coverage in brackets. The last three columns give the evidence that would favour the classical method ($\to$C), favour a quantum workflow ($\to$Q), or firm up the resource estimate ($\to$R). Each is a question that would change some assessment, not a sufficient condition for moving a whole class. Letter codes as in Fig.~\ref{fig:taxonomy}.}\label{tab:forward}\\
\toprule
\# & Class & Present status & Resource evidence [coverage] & $\to$C & $\to$Q & $\to$R\\
\midrule\endfirsthead
\toprule
\# & Class & Present status & Resource evidence [coverage] & $\to$C & $\to$Q & $\to$R\\
\midrule\endhead
\bottomrule\endfoot
1 & OTOC(2) scrambling, 2D & \psQ{} (\eld{E06} d), conditional & 65\,q measured; no FT compilation [\mtN] & Converged classical reproduction of the stated observable, depth and precision & Validated accuracy and full cost at larger scale & FT compilation\\
2 & Operator Loschmidt echo & \psQ, conditional & 56\,q measured [\mtN] & Majorana/Pauli propagation converged over the full regime & Bounded-error estimator costed for the same regime & FT compilation\\
3 & Floquet prethermal 2D magnetism: H2 and Heron instances & \psQ{} (\eld{E05} d; \eld{E23} d), two separate instances & 56\,q and 74\,q measured [\mtN] & Converged TN, NQS or Pauli-path result in the hardest regime of either instance & Validation across the hard regime, mitigation overhead costed & FT compilation\\
4 & 2D TFIM dynamics, 100 spins, 10 Trotter steps & \psN & 230 LQ, $2.4\times10^6$ T per execution [\mtP] & Converged classical result for the stated time, state and observable & Costed estimation below classical & Repetitions, preparation\\
5 & 3D spin-glass quenches & \psC{} (\mech{S}, \mech{V}), analog scope & Analog only [\mtN] & Stable long-time t-VMC or 3D BP-TN & Observables beyond classical coverage & Gate-model compilation\\
6 & 2D frustrated $J_1$--$J_2$ Heisenberg, $10\times10$ cylinder & \psN{} (baseline is 4-core DMRG) & $5.6\times10^4$\,s single factory [\mtP] & NQS, PEPS or GPU DMRG on the same cylinder at $\epsilon=0.01$ total, with measured cost & Validated adiabatic preparation & Gap scaling\\
7 & 2D Fermi--Hubbard energy, half filling & \psN{} (\mech{S}: sign-free; best classical baseline unresolved) & $1.05\times10^5$\,s at $10\times10$ vs 4-core DMRG $7.1\times10^9$\,s [\mtP] & QMC or parallel DMRG at $\epsilon=0.01$ total with measured time below the modelled schedule & Estimate below the best measured classical cost & Preparation\\
8 & Strongly correlated materials & \psN & No multi-orbital estimate [\mtN] & Downfolded classical solution at decision accuracy & Multi-orbital estimate with downfolding error and decision tolerance & Multi-orbital estimate\\
9 & Core-level XAS & \psN{} (resource-estimated candidate) & \xasCirc{} Toffoli longest circuit; \xasAlg{} total; \xasHours\,h at 1\,MHz [\mtC, assumptions unvalidated] & Classical spectrum at decision accuracy from the best available method & Validated preparation and error model, cost below that spectrum & Physical-level compilation\\
10 & FeMoco precision and other models & \psN & Dense code: \femoT\,h expected at \femoSevMinvolQ\,M physical, \femoTofMinvol{} Toffoli; THC $3.2\times10^{10}$ Toffoli [\mtP] & Classical estimate at matched confidence for larger models & Measured overlap; complete protocol & Larger-model counts\\
11 & P450 Cpd I spin gaps & \psN & 2022: $7.8\times10^9$ Toffoli, 73\,h; 2026: \pfMinvolH\,h at 1.03\,M physical [\mtP] & Converged classical spin gap at a stated tolerance & Cost for both spin states with error cancellation & Preparation, success criterion and both spin states on an architecture\\
12 & Ru CO$_2$-fixation catalysts & \psN & 2740 LQ, $5.4\times10^{11}$ T (older model); 56-orbital model 0.46\,h at 1.57\,M physical [\mtP] & Classical multireference result & Model-specific modern count & Modern count, older model\\
13 & Transition-metal ML labels & \psN & --- [\mtN] & Classical labels adequate & Incremental value of better labels against the downstream error budget & Label cost\\
14 & Enzyme active sites & \psN & No estimate for a named cluster [\mtN] & Embedding with a classical or neural solver adequate for a named cluster and observable & Solver error dominating a validated budget & Embedded estimates\\
15 & Learning from quantum data & \psT{} (quantum-side sample-complexity separation) & Separation with quantum memory~\cite{huang2022exp} [n/a] & Equivalent classical access model & Task demonstrations at scale & Copies, memory, access\\
16 & Real-time SU(3) gauge theory & \psN & $10^9$--$6.5\times10^{48}$ T across formulations, evolution only [\mtS] & Classical real-time method for a controlled physical observable & Controlled-truncation estimate incl.\ preparation and measurement & Single formulation, full observable\\
17 & Sampling capability: RCS, GBS, doped Clifford, analog atoms & \psCap & --- & --- & --- & ---\\
18 & SQD, single-layer circuit-energy objective & \psM{} (\mech{B}); workflow as a whole \psN & --- & --- & Sampler ablation under a common selection, recovery and diagonalization procedure & ---\\
19 & FeMoco-76 model energy & \psV{} (feasibility, extrapolation-dependent; operational ranking unresolved) & --- & --- & --- (a failed estimate would reopen the comparison) & ---\\
20 & P450 $\sim$500-orbital model & \psN & 9000 LQ, $1.5\times10^{12}$ Toffoli [\mtP] & Embedding solves the decision & A gap or decision target with a complete comparison & Modern count\\
21 & Generic ground-state chemistry & \psN & --- [\mtN] & --- & Representative family, fixed accuracy, with empirical or proven exponential classical scaling & ---\\
22 & Main-group thermochemistry & \psV{} (named benchmarks; \mech{R}) & --- & --- & Validated-domain failure with consequence (a solver need, not yet a quantum result) & ---\\
23 & Rydberg topological diagnostics & \psN & --- [\mtN] & --- & A matched computational task & Measurement cost\\
24 & High-throughput materials discovery & \psV{} (\mech{A}; experimentally validated~\cite{chen2024}) & --- & --- & Decision-critical electronic-structure bottleneck identified & Bottleneck estimate\\
25 & Binding free energies, MD & \psV{} (mature practice; \mech{A}); one unmatched hybrid workflow~\cite{merz2026} & --- & --- & Quantum subproblem dominating the budget, shown by ablation & Subproblem estimate\\
26 & Property learning within a phase & \psT{} (classical-side guarantee in theorem scope; \mech{A}) & --- & --- & Labels unobtainable classically & ---\\
27 & Protein structure prediction, proteome scale & \psV{} (\mech{R}); fragments are a separate record & --- & --- & Fragment gain under identical selection and cost rules & ---\\
28 & Generative molecular design & \psN & --- [\mtN] & Matched ablation shows no gain & Matched ablation incl.\ downstream selection shows gain & ---\\
29 & Unstructured search & \psT{} (conditional model; query-model limit~\cite{turturean2026}; \mech{H}) & Model only [\mtS] & --- & $t_Q\lesssim10\,\mu$s with $Q\gg1$ within a stated latency, or a structured speedup & Oracle estimate\\
30 & Approximate optimization: annealing (QAC) and Rydberg MIS & \psC, family-specific & QAC $>$1300 encoded variables; Rydberg $\le$289\,q [\mtN] & Classical portfolio matches scaling on the defined family & Separation vs the full portfolio with stated runtime accounting & ---\\
31 & QML on classical data & \psN & Constructed separations~\cite{liu2021} [\mtN] & --- & Practical separation on a stated workload and access model & ---\\
32 & Nuclear ground states, $A\le20$ & \psV{} (cited interactions; \mech{V}); no quantum comparator identified & --- & --- & A quantum comparator in scope at matched precision & ---\\
33 & Euclidean lattice QCD, zero density & \psV{} (mature practice); no demonstrated practical matched quantum advantage & Small-model sampling proposal~\cite{yamamoto2022} [\mtN] & --- & Lower autocorrelation and full observable cost than modern HMC at relevant lattice spacing & ---\\
34 & CFD: full field / selected outputs & \psV{} (full field: no costed quantum workflow) / selected outputs: see class 38 & Bounded advantage, modest for selected observables~\cite{jennings2025} [\mtS] & --- & Selected-observable workload with value and an end-to-end cost below explicit classical CFD & Selected-output estimate at an engineering Reynolds number\\
35 & Weather/climate: forecasts / statistics & \psV{} per named forecast or component / \psN & --- [\mtN] & --- & A named component with a costed advantage & Component estimate\\
36 & Doped 2D Fermi--Hubbard (sign-problem instances) & \psN{} (new record) & --- [\mtN] & NQS, PEPS, constrained-path or embedding result at stated density, $U/t$, $t'$ and tolerance, with measured cost & Costed QPE estimate for the same instance & Doped-lattice estimate\\
37 & Quantum-enhanced MCMC sampling & \psN{} (conditional candidate) & Proposals and early hardware experiments~\cite{layden2023,marshall2026} [\mtN] & Classical nonlocal updates matching time per independent sample & Lower time per independent sample incl.\ proposal generation and noise & Proposal circuit cost\\
38 & Sparse linear solves and time stepping for engineering PDEs (FEM, CFD) & \psT{} (at most polynomial speedup in fixed dimension~\cite{montanaro2016}; \mech{IO}, \mech{H}) & $10^{21}$--$10^{39}$ logical qubits $\times$ T gates for drag on a sphere~\cite{penuel2025}; depth $1.8\times10^{29}$ for EM scattering~\cite{scherer2017}; one crossover at $2^{80}$ grid points~\cite{zhuang2025} [\mtP] & --- (classical practice is the incumbent) & A high-dimensional or quantum-native linear problem with scalar output and an end-to-end cost below tuned multigrid or Monte Carlo & Constant-factor estimate at industrial condition numbers with input and readout\\
\end{longtable}
}


% =====================================================================
\section{Resource models, records and further figures}
\label{app:models}

\paragraph{Rows shown in Fig.~\ref{fig:dc76}.} The rows are from Table~XXVII of the dense-code compilation (physical error rate $10^{-3}$, 1\,$\mu$s code cycle, 10\,$\mu$s reaction time). From left to right: the minimum-qubit row, the rows the source highlights as running within a month and within a day, the minimum-volume row and the fastest tabulated row; for FeMoco a second fast row at \femoSevSelQ\,M qubits is also shown. The vertical axis is the source's expected runtime, which is the one-shot runtime divided by one minus the logical failure probability (its Sec.~VI.B). At the low-qubit end most shots fail: the minimum-qubit FeMoco row has a failure probability of \femoMinqPfail, so \femoMinqShot{} hours per shot become \femoMinqH{} expected hours. Spectral-hit repetitions, state preparation and the second spin state needed for a P450 gap are not included, and classical reference values are given in the text with their own scopes.

\subsection{Runtime bounds and implied rates}
In general a fault-tolerant execution is bounded below by its slowest required resource stream,
\begin{equation}
T_{\rm run}\ \gtrsim\ \max\!\left\{\frac{N_T}{R_T},\ \frac{N_{\rm CCZ}}{R_{\rm CCZ}},\ D_{\rm adapt}\,t_{\rm react},\ D_{\rm code}\,t_{\rm cycle},\ T_{\rm routing/readout}\right\},
\label{eq:bottleneck}
\end{equation}
where $R_T$ and $R_{\rm CCZ}$ are the production rates of T and CCZ states, $D_{\rm adapt}$ the number of sequential reaction-dependent stages and $D_{\rm code}$ the logical depth in code cycles. State preparation, resets and outer repetitions come on top. Inverting published estimates gives the effective rates of Fig.~\ref{fig:rates}. A complete execution record would list
$(Q_{\rm data}, Q_{\rm work}, Q_{\rm factory}, N_{\rm T}, N_{\rm CCZ}, D_{\rm adapt}, p_{\rm fail}, t_{\rm cycle}, t_{\rm react})$; most published estimates report a subset.

\begin{figure}[htbp]
\centering
\begin{tikzpicture}
\begin{axis}[
  xbar, xmode=log, log origin=infty, width=0.66\textwidth, height=6.6cm,
  xmin=10, xmax=2e7, y dir=reverse, ytick={1,...,8}, ymin=0.4, ymax=8.6,
  yticklabels={{TFIM, ns gates, $p{=}10^{-4}$}, {TFIM, ns gates, $p{=}10^{-3}$}, {Ru catalyst, ns, $p{=}10^{-4}$}, {Ru catalyst, ns, $p{=}10^{-3}$}, {TFIM, $\mu$s gates, $p{=}10^{-4}$}, {Ru catalyst, $\mu$s, $p{=}10^{-4}$}, {P450 model X, THC, $p{=}10^{-3}$}, {FeMoco-54, THC, 4\,M phys.}},
  yticklabel style={font=\footnotesize},
  xlabel={\footnotesize Effective non-Clifford rate $N/T$ (states per second)},
  xticklabel style={font=\footnotesize}, bar width=7pt, bar shift=0pt,
  xmajorgrids, grid style={black!10},
  nodes near coords={\pgfmathprintnumber[sci,precision=1]{\pgfplotspointmeta}}, point meta=x, every node near coord/.append style={font=\scriptsize, anchor=west},
  axis x line*=bottom, axis y line*=left,
  legend style={font=\scriptsize, at={(0.98,0.03)}, anchor=south east, draw=black!25, fill=white}]
\addplot[fill=qcF, draw=none] table[x=x,y=y] {figdata/rates_T.dat};
\addlegendentry{T states (Beverland et al.)}
\addplot[fill=black!20, draw=black!75, line width=0.3pt, postaction={pattern=north east lines, pattern color=black!80}] table[x=x,y=y] {figdata/rates_CCZ.dat};
\addlegendentry{Toffoli/CCZ}
\end{axis}
\end{tikzpicture}
\caption[Rates implied by particular schedules]{Effective non-Clifford rates $N/T$ implied by particular published schedules, from \texttt{resources.csv}. T-state and Toffoli rates are different quantities and are drawn with different fills (plain bars for T states, hatched bars for Toffoli gates); the validator refuses a Toffoli record in the T series. Each value belongs to its source schedule (factory count, code, width, gate speed) and is not a portable platform performance. Sources: TFIM and Ru catalyst~\cite{beverland2022}; P450, $7.8\times10^9$ Toffoli in 73\,h~\cite{goings2022}; FeMoco 54-orbital layout, $6.7\times10^9$ Toffoli in 3 days at a stated 25\,kHz~\cite{lee2021thc}.}
\label{fig:rates}
\end{figure}


\subsection{Comparison with announced hardware targets}
\label{app:vendor}
Public hardware roadmaps provide scale references, not application schedules (Table~\ref{tab:vendor}). T counts, Toffoli counts and vendor ``operations'' are different units, and the vendor unit is not defined in the announcements. We therefore compare only under a named accounting convention and draw only one-sided conclusions about named compiled circuits.

\begin{rulebox}{Convention X and bound logic}
\textbf{Convention X (illustrative):} one vendor operation implements at most one counted non-Clifford gate (T or CCZ alike), with Clifford, routing and memory operations free. This is an accounting convention, not a property of any architecture. A source circuit's count is not a lower bound on every circuit for the same task.
\begin{itemize}
\item If a \emph{named compiled circuit} has count $N\ge L$ and $L>B$, that circuit does not fit budget $B$ under convention X or any costlier one: a count obstruction for that circuit.
\item If $N\le U$ and $U\le B$, the count is within $B$ under convention X only; compatibility with a real architecture is \emph{not established}.
\item A workload of many independent circuits is tested one circuit at a time. A cumulative count is comparable only with a cumulative budget, which the announcements do not state.
\item A compiled width above the announced logical-qubit count, under the same width convention, is a width obstruction as compiled. A record with no reported width has unknown compatibility.
\end{itemize}
\end{rulebox}
Logical error and circuit volume must also agree. Under an independent-failure model with failure probability $p_L$ per counted operation, $P(\text{success})\approx e^{-Gp_L}$, so $G=10^8$ operations at 99\% success need $p_L\lesssim10^{-10}$ and $G=10^9$ need $p_L\lesssim10^{-11}$.

\begin{table}[htbp]
\centering\small
\caption{Announced vendor targets used as scale references. Targets are forward-looking statements, not delivery commitments.}\label{tab:vendor}
\begin{tabular}{@{}L{3.1cm} L{1.6cm} L{2.4cm} L{2.3cm} L{5.0cm}@{}}
\toprule
System & Target date & Logical qubits & Logical operations & Error target and implied success\\\midrule
IBM Starling~\cite{ibmroadmap2025} & 2029 & 200 & $10^8$ & Unit not specified in announcement; $10^8$ ops at 99\% success implies $p_L\lesssim10^{-10}$\\
IBM Blue Jay~\cite{ibmroadmap2025,hpcwire2025} & 2033+ & 2000 & $10^9$ & As above; implies $p_L\lesssim10^{-11}$\\
Quantinuum Apollo~\cite{quantinuum2024,quantFiling} & 2029 & hundreds & not stated as an operation count & Logical fidelity target ``Ten Nines'', error as low as $\sim10^{-10}$ per filing\\
\bottomrule
\end{tabular}
\end{table}

\begin{figure}[htbp]
\centering
\begin{tikzpicture}
\begin{axis}[
  xmode=log, ymode=log, width=0.82\textwidth, height=8.2cm,
  xmin=60, xmax=1e5, ymin=1e6, ymax=1e13,
  xlabel={\footnotesize Logical qubits, in each source's width convention}, ylabel={\footnotesize Non-Clifford count in source unit},
  ticklabel style={font=\footnotesize}, grid=major, grid style={black!10}, clip=false,
  legend style={font=\scriptsize, at={(0.5,-0.16)}, anchor=north, draw=black!25, fill=white, legend columns=2}, legend cell align=left]
\draw[dashed, draw=qcQ, line width=0.8pt] (axis cs:60,1e8) -- (axis cs:200,1e8) -- (axis cs:200,1e6);
\draw[dashed, draw=qcQ, line width=0.8pt] (axis cs:60,1e9) -- (axis cs:2000,1e9) -- (axis cs:2000,1e6);
\node[font=\scriptsize, qcQ, anchor=west] at (axis cs:62,2.5e7) {200 LQ / $10^8$ ops};
\node[font=\scriptsize, qcQ, anchor=south west] at (axis cs:62,1.15e9) {2000 LQ / $10^9$ ops (IBM targets, per circuit)};
\addplot[only marks, mark=*, mark size=2.6pt, color=qcF] table[x=x,y=y] {figdata/resmap_T_point.dat};
\addlegendentry{T count, one circuit}
\addplot[only marks, mark=triangle*, mark size=3.2pt, color=black!70] table[x=x,y=y] {figdata/resmap_Toff_point.dat};
\addlegendentry{Toffoli count, one circuit}
% generated by ledger/make_figdata.py
\node[font=\scriptsize, inner sep=5pt, anchor=west] at (axis cs:230,2.4e6) {2D TFIM dynamics, per execution};
\node[font=\scriptsize, inner sep=5pt, anchor=north] at (axis cs:2740,5.4e11) {Ru CO$_2$ catalyst};
\node[font=\scriptsize, inner sep=5pt, anchor=west] at (axis cs:25481,1.5e10) {RSA-2048};
\node[font=\scriptsize, inner sep=5pt, anchor=west] at (axis cs:100,3.54e8) {XAS, longest circuit (min.\ workspace)};
\node[font=\scriptsize, inner sep=5pt, anchor=north west] at (axis cs:2142,5.3e9) {FeMoco-54 (Reiher), THC};
\node[font=\scriptsize, inner sep=5pt, anchor=west] at (axis cs:2196,3.2e10) {FeMoco-76 (Li), THC};
\node[font=\scriptsize, inner sep=5pt, anchor=south east] at (axis cs:9000,1.5e12) {P450 $\sim$500-orbital model};

\end{axis}
\end{tikzpicture}
\caption[Single-circuit resource estimates and announced vendor targets]{Published single-circuit resource estimates in their original units; only records read at table or section level with a reported compiled width are shown. Dashed lines are announced IBM targets, drawn as reference lines and not as feasibility regions. By the bound logic of Appendix~\ref{app:vendor}, a point beyond a line is an obstruction for that compiled circuit under convention X; a point inside a line establishes nothing about a real architecture. The XAS point is the longest circuit at the minimum algorithm workspace of 100 logical qubits, excluding factories~\cite{fomichev2025}. Cumulative counts over many circuits, such as the \xasAlg{} Toffoli gates of the full XAS spectrum, are not comparable with a per-circuit budget and are not drawn. The two FeMoco points are different Hamiltonians~\cite{lee2021thc}. Records without a reported width (lattice models, dense-code compilations, most gauge formulations) are omitted. Sources: TFIM, Ru, RSA~\cite{beverland2022}; P450~\cite{goings2022}.}
\label{fig:resource}
\end{figure}


\subsection{Break-even and output scenarios}
Fig.~\ref{fig:quadratic} plots Eq.~(\ref{eq:tstar}). The cubic and quartic curves are sensitivities for hypothetical structured algorithms (\beCubic\,s and \beQuartic\,s at $t_Q=1$\,ms and $P=10^6$); they are not known search algorithms.

\begin{figure}[htbp]
\centering
\begin{tikzpicture}
\begin{axis}[
  xmode=log, ymode=log, width=0.56\textwidth, height=6.4cm,
  xmin=1, xmax=1e7, ymin=1e-3, ymax=1e10,
  xlabel={\footnotesize Effective classical parallelism $P$ (cores)},
  ylabel={\footnotesize Break-even runtime $T^*$ (s)},
  ticklabel style={font=\footnotesize}, grid=major, grid style={black!10},
  legend style={font=\scriptsize, at={(1.04,0.5)}, anchor=west, draw=black!25, fill=white}, legend cell align=left]
\addplot[qcF, line width=1.3pt, mark=*, mark size=1.8pt, mark repeat=5] table[x=P,y=d2ms] {figdata/breakeven_curves.dat};
\addlegendentry{$d=2$, $t_Q=1$\,ms, $Q=1$}
\addplot[qcF!75, line width=1.0pt, mark=square*, mark size=1.7pt, mark repeat=5, mark phase=2] table[x=P,y=d2hundred] {figdata/breakeven_curves.dat};
\addlegendentry{$d=2$, $t_Q=100\,\mu$s, $Q=1$ (= 1\,ms, $Q=100$)}
\addplot[qcF!55, line width=1.0pt, mark=triangle*, mark size=2.2pt, mark repeat=5, mark phase=4] table[x=P,y=d2ten] {figdata/breakeven_curves.dat};
\addlegendentry{$d=2$, $t_Q=10\,\mu$s, $Q=1$}
\addplot[black!60, line width=1.0pt, dashed] table[x=P,y=d3ms] {figdata/breakeven_curves.dat};
\addlegendentry{hypothetical $d=3$, $t_Q=1$\,ms}
\addplot[black!60, line width=1.0pt, dotted] table[x=P,y=d4ms] {figdata/breakeven_curves.dat};
\addlegendentry{hypothetical $d=4$, $t_Q=1$\,ms}
\draw[black!40, dashed] (axis cs:1,86400) -- (axis cs:1e7,86400) node[pos=0.02, above, font=\scriptsize, anchor=south west]{1 day};
\draw[black!40, dashed] (axis cs:1,3.16e7) -- (axis cs:1e7,3.16e7) node[pos=0.02, above, font=\scriptsize, anchor=south west]{1 year};
\node[circle, fill=qcF, inner sep=1.5pt] at (axis cs:1e6,1e8) {};
\node[font=\scriptsize, anchor=east, inner sep=5pt] at (axis cs:1e6,1e8) {3.2 y};
\node[circle, fill=qcF!70, inner sep=1.5pt] at (axis cs:1e6,1e6) {};
\node[font=\scriptsize, anchor=west, inner sep=5pt] at (axis cs:1e6,1e6) {11.6 d};
\node[circle, fill=qcF!45, inner sep=1.5pt] at (axis cs:1e6,10000) {};
\node[font=\scriptsize, anchor=west, inner sep=5pt] at (axis cs:1e6,10000) {2.8 h};
\node[circle, fill=black!60, inner sep=1.5pt] at (axis cs:1e6,4.64159) {};
\node[font=\scriptsize, anchor=north, inner sep=5pt] at (axis cs:1e6,4.64159) {4.6 s};

\end{axis}
\end{tikzpicture}
\caption[Break-even runtime for polynomial speedups]{Break-even runtime $T^*$ (Eq.~\ref{eq:tstar}) with $t_C=10$\,ns per equivalent oracle evaluation; a formula figure whose parameters are in \texttt{models.csv}. Solid lines with markers are the quadratic case (circles, squares and triangles mark the three gate times of the legend), which corresponds to unstructured search. Dashed and dotted lines are model sensitivities for hypothetical structured algorithms of higher degree, not known search algorithms. The model assumes equivalent oracle work on both sides, effective parallelism $P$, a problem that partitions across $Q$ quantum processors, and reaction-dependent stages that set $t_Q$. No acceptable-latency threshold is implied.}
\label{fig:quadratic}
\end{figure}


For output, suppose a receiver needs $N$ values of $b$ bits that are independent of everything it already holds, obtained from fresh preparations of an $n$-qubit state. Each shot yields at most $n$ classical bits, so
\begin{equation}
T^{\rm out}_Q \;\ge\; \frac{N\,b\,t_{\rm shot}}{n\,Q},\qquad T^{\rm out}_C \;\approx\; \frac{N\,b}{B_{\rm mem}},
\label{eq:output}
\end{equation}
for $Q$ quantum devices and a classical memory bandwidth $B_{\rm mem}$. With $N=10^9$, $b=32$, $n=10^3$, $t_{\rm shot}=1$\,ms and $B_{\rm mem}=10^{12}$\,bits/s the floor is \outQ\,s per quantum device against \outC\,s of classical writing (Fig.~\ref{fig:readout}). The scenario is a conditional floor for one information model. A simulated field is a deterministic function of its inputs, so it need not carry $Nb$ bits of new information; a smooth field has a short description; and an algorithm that outputs selected observables is not subject to the floor.

\begin{figure}[htbp]
\centering
\begin{tikzpicture}
\begin{axis}[
  xmode=log, ymode=log, width=0.66\textwidth, height=6.0cm,
  xmin=1e6, xmax=1e21, ymin=1e-6, ymax=1e18,
  xlabel={\footnotesize Independent information to deliver: $N$ values of 32 bits},
  ylabel={\footnotesize Output time (s)},
  ticklabel style={font=\footnotesize}, grid=major, grid style={black!10},
  legend style={font=\scriptsize, at={(0.02,0.98)}, anchor=north west, draw=black!25, fill=white}, legend cell align=left]
\addplot[qcF, line width=1.2pt] table[x=N,y=q_ms] {figdata/readout_curves.dat};
\addlegendentry{quantum floor, $n=10^3$, $t_{\rm shot}=1$\,ms}
\addplot[qcF!55, line width=1.1pt, dashed] table[x=N,y=q_us] {figdata/readout_curves.dat};
\addlegendentry{quantum floor, $n=10^3$, $t_{\rm shot}=1\,\mu$s}
\addplot[black!75, line width=1.2pt, dashdotted] table[x=N,y=c_node] {figdata/readout_curves.dat};
\addlegendentry{classical node, $10^{12}$ bits/s}
\draw[black!40, dashed] (axis cs:1e6,86400) -- (axis cs:1e21,86400) node[pos=0.98, above, anchor=south east, font=\scriptsize]{1 day};
\draw[black!40, dashed] (axis cs:1e6,3.16e7) -- (axis cs:1e21,3.16e7) node[pos=0.98, above, anchor=south east, font=\scriptsize]{1 year};
\end{axis}
\end{tikzpicture}
\caption[Output time under an information-rate scenario]{An information-rate scenario for output (Eq.~\ref{eq:output}): time to deliver $N$ independent 32-bit values from one quantum device by fresh preparations of a 1{,}000-qubit state, against writing them on one classical node. The scenario assumes that the receiver holds no information that shortens the description, which a deterministic simulated field need not satisfy. The curves exclude the solve on both sides, do not rank total runtimes and do not apply to algorithms that output selected observables. A formula figure; parameters in \texttt{models.csv}.}
\label{fig:readout}
\end{figure}


\subsection{Resource records}
{\scriptsize
\setlength{\tabcolsep}{3pt}
\begin{longtable}{@{}L{2.5cm} L{1.15cm} L{1.55cm} L{2.85cm} L{3.0cm} L{3.05cm} L{0.95cm}@{}}
\caption{Resource records used in this paper, one per model, algorithm and schedule (\texttt{resources.csv}). LQ = logical qubits as compiled unless stated. ``Capability'' applies the bound logic of Appendix~\ref{app:vendor} under convention X to the named circuit against the IBM targets (200 LQ / $10^8$; 2000 LQ / $10^9$). Coverage: \mtC{} = preparation, observable, repetitions and error budget costed for the stated protocol; \mtP{} = listed omissions; \mtS{} = evolution-only or truncation-dependent. Coverage says nothing about whether the assumptions are validated. Rows marked $\ddagger$ carry a value known at second hand or from an abstract only; no such value is used in a figure or ratio. $^\S$n.c.\ = not constructed: the dense-code source tabulates compute, hot-storage and cold-storage regions separately, and we do not add them into a logical width.}\label{tab:worked}\\
\toprule
Record & LQ & $N_{\rm nc}$ & Quantum time & Classical time & Capability (convention X) & Cov.\\
\midrule\endfirsthead
\toprule
Record & LQ & $N_{\rm nc}$ & Quantum time & Classical time & Capability (convention X) & Cov.\\
\midrule\endhead
\bottomrule\endfoot
2D TFIM dynamics, 100 spins, 10 steps~\cite{beverland2022} & 230 & $2.4\times10^6$ T & 0.56\,s per execution (ns, $p{=}10^{-4}$); repetitions excluded & No current baseline compiled & Width obstruction vs 200 as compiled; count within $10^8$ under X only & \mtP\\
Fermi--Hubbard $6\times6$ cylinder, $U/t{=}4$, half filling, $\epsilon=0.01$ total~\cite{yoshioka2024} & 72 (system) & not reported & $1.23\times10^4$\,s (1 factory, $4.3\times10^5$ physical); $1.59\times10^3$\,s (16 factories, $3.4\times10^6$) & DMRG $2.2\times10^5$\,s on 4 cores (est.); QMC sign-free, not costed & --- & \mtP\\
Fermi--Hubbard $10\times10$, same~\cite{yoshioka2024} & 200 (system) & not reported & $1.05\times10^5$\,s (1 factory, $8.9\times10^5$ physical); $1.46\times10^4$\,s (16 factories, $5.4\times10^6$) & DMRG $7.1\times10^9$\,s on 4 cores (est.) = \hubTenCoreH{} core-h; QMC not costed & --- & \mtP\\
$J_1$--$J_2$ Heisenberg $10\times10$ cylinder, $J_2/J_1{=}0.5$~\cite{yoshioka2024} & 100 (system) & not reported & $5.64\times10^4$\,s (1 factory, $5.1\times10^5$ physical); $4.32\times10^3$\,s (16 factories, $4.4\times10^6$) & DMRG $7.9\times10^4$\,s on 4 cores (est.); NQS on another geometry: 4 d on 20 GPUs~\cite{rende2024} & --- & \mtP\\
XAS, Li$_4$Mn$_2$O, CAS(22e,18o)~\cite{fomichev2025}: longest circuit & 100 (min.\ work\-space, excl.\ factories) & \xasCirc{} Toffoli & \xasCircS\,s at a 1\,MHz logical clock ($1.08\times10^8$ cycles) & --- & Width within 200 by that convention; count obstruction vs $10^8$ for this circuit; within $10^9$ under X only & \mtC\\
XAS, same: full sampled spectrum & 350 (active-volume implementation) & \xasAlg{} Toffoli over all circuits & \xasHours\,h at 1\,MHz (source: less than two days) & RAS: $\sim$32 days for 1{,}000 states (source's estimate, which it calls likely optimistic on the classical side) & Cumulative count; not comparable with a per-circuit budget & \mtC\\
FeMoco, 76 orbitals, THC~\cite{lee2021thc} & 2196 & $3.2\times10^{10}$ Toffoli & no physical layout in the source & --- & Count obstruction vs $10^9$ for this algorithm & \mtP\\
FeMoco, 54 orbitals, THC~\cite{lee2021thc} & 2142 & $5.3\times10^9$ Toffoli ($6.7\times10^9$ in the layout) & 3 days at 25\,kHz on 4\,M physical (abstract: under four days) & --- & Count obstruction vs $10^9$ for this algorithm & \mtP\\
FeMoco, 76 orbitals, dense code, minimum volume~\cite{low2026dense} & n.c.$^\S$ & \femoTofMinvol{} Toffoli & \femoT\,h expected (\femoShot\,h per shot, $p_{\rm fail}=\femoPfail$), \femoSevMinvolQ\,M physical & Three classical records (Section~\ref{sec:chem}) & Count obstruction vs $10^9$ for this compilation & \mtP\\
FeMoco, 76 orbitals, dense code, fastest tabulated~\cite{low2026dense} & n.c.$^\S$ & \femoTofFast{} Toffoli & \femoFastT\,h expected, \femoSevFastQ\,M physical & as above & Count obstruction vs $10^9$ & \mtP\\
FeMoco, 76 orbitals, dense code, within a month~\cite{low2026dense} & n.c.$^\S$ & $1.49\times10^{10}$ Toffoli & 716\,h expected (529\,h per shot, $p_{\rm fail}=\femoLowPfail$), \femoSevLowQ{} physical & as above & Count obstruction vs $10^9$ & \mtP\\
FeMoco, 76 orbitals, spectrum amplification (count)~\cite{low2025sa} & 1459 & $9.99\times10^8$ Toffoli per run & no physical layout with this count & --- & Width within 2000; count within $10^9$ under X only, with no margin & \mtP\\
FeMoco, 54 orbitals, spectrum amplification (count)~\cite{low2025sa} & 1137 & $3.41\times10^8$ Toffoli per run & no physical layout with this count & --- & Width within 2000; count within $10^9$ under X only & \mtP\\
FeMoco, 76 orbitals, spectrum amplification (schedule)$^\ddagger$~\cite{low2025sa,zhai2026} & not extracted & not extracted & 8.6\,h per attempt, four CCZ factories, $\sim$4.5\,M physical & --- & --- & \mtP\\
FeMoco, candidate-refinement protocol~\cite{berry2024} & not extracted & $7.3\times10^{10}$ Toffoli (2 or 3 phase-estimation rounds; 2.3 times the single-shot THC count) & --- & --- & Count obstruction vs $10^9$ & \mtP\\
FeMoco, 76 orbitals, classical campaign~\cite{zhai2026} & --- & --- & --- & 2.77\,M core-h measured; $\pm$0.31/$\pm$1.25 kcal/mol & --- & ---\\
FeMoco, 54 orbitals, dense code, minimum volume~\cite{low2026dense} & n.c.$^\S$ & $4.55\times10^8$ Toffoli & 1.13\,h expected, 1.15\,M physical (Fig.~\ref{fig:femoco}) & --- & Within $10^9$ under X only & \mtP\\
P450 Cpd I model X, 2022~\cite{goings2022} & not reported$^\dagger$ & $7.8\times10^9$ Toffoli & 73\,h on 4.6\,M physical ($p{=}10^{-3}$); 25\,h on 0.5\,M ($10^{-5}$) & DMRG $M{=}3000$: 36{,}564 CPU-h, $\sim$1 month (both estimated in source) & Count obstruction vs $10^9$ for this algorithm & \mtP\\
P450 Cpd I model X (58o63e), dense code, 2026~\cite{low2026dense} & n.c.$^\S$ & \pfTofMinvol{} (minimum volume); \pfTofFast{} (fastest) Toffoli & \pfMinvolH\,h at 1.03\,M physical; \pfFastH\,h at 2.18\,M; \pfMinqDays{} days at 81{,}000 & as above & Within $10^9$ under X only & \mtP\\
P450 $\sim$500-orbital model~\cite{goings2022} & 9000 & $1.5\times10^{12}$ Toffoli & --- & --- & Width and count obstructions & \mtP\\
Ru CO$_2$-fixation catalyst~\cite{beverland2022} & 2740 & $5.4\times10^{11}$ T & 1 month (ns); 130\,y ($\mu$s) & Not reported & Width and count obstructions for this circuit & \mtP\\
Ru catalyst XVIII, 56o64e (another active space), dense code~\cite{low2026dense} & n.c.$^\S$ & $2.25\times10^8$ Toffoli (fastest row); $5.55\times10^9$ (minimum qubits) & 0.46\,h at 1.57\,M physical; 3{,}200\,h at 76{,}000 & Not reported & --- & \mtP\\
EM scattering by finite elements, $3.3\times10^8$ unknowns~\cite{scherer2017} & $\sim10^8$ (width incl.\ oracles) & depth $1.8\times10^{29}$; T-depth $8.2\times10^{28}$ & --- & big-O comparison only & Width and count obstructions & \mtP\\
Drag on a sphere, Carleman lattice Boltzmann$^\ddagger$~\cite{penuel2025} & \multicolumn{2}{L{3.0cm}}{$10^{21}$--$10^{39}$ logical qubits $\times$ T gates} & --- & steady RANS on one workstation, 1--10\,h for $\mathrm{Re}=10^3$--$10^7$ (source's reference); explicit CFD scales as $O(\mathrm{Re}^3)$ & not assessed (abstract-level value) & \mtP\\
Navier--Stokes, $2^{80}$-point mesh~\cite{zhuang2025} & 181 & $9.41\times10^7$ Toffoli and $3.94\times10^8$ rotations per preparation; 1{,}000 preparations & 42.6 days, 8.71\,M physical (94\% factories), distance 25 & 130 years: banded-Cholesky operation count for the full system at a benchmark rate & Width within 200; count within $10^8$ per preparation under X only & \mtP\\
Drag on a sphere, public per-instance records~\cite{qsub2025}: $\mathrm{Re}=10$ and $10^8$ & 5{,}679; $8.6\times10^{10}$ & $9.3\times10^{17}$; $9.2\times10^{29}$ T & --- & as above & not assessed (unreconciled with the published range) & \mtP\\
SU(3) shear viscosity, $\Sigma(72\times3)$ subgroup~\cite{osorio2025} & 27{,}000 (data registers only) & $3.5\times10^{12}$ T (Kogut--Susskind); $7.1\times10^{12}$ T (improved) & evolution only & as below & Width and count obstructions & \mtS\\
SU(3) plaquette block encoding, one oracle~\cite{draper2026} & not a total width & $1.2\times10^5$ T ($B{=}6$); $1.2\times10^6$ T ($B{=}9$) & subroutine only & --- & not an application count & \mtS\\
SU(3) shear viscosity, $10^3$: separate formulations (Fig.~\ref{fig:lgt})~\cite{lamm2026,ciavarella2025} & per formulation; mostly not reported & $10^9$ to $6.5\times10^{48}$ T & evolution only & No scalable classical route known for the full theory & The one record with a width does not fit; unknown where no width is reported & \mtS\\
Unstructured search~\cite{babbush2021,hoefler2023} & any & $\sqrt{M/Q}$ calls & Eq.~(\ref{eq:tstar}) & $Mt_C/P$ & --- & \mtS\\
RSA-2048 (reference)~\cite{beverland2022} & 25{,}481 & $1.5\times10^{10}$ T & 18\,h ($10^{-4}$, ns) & --- & Width and count obstructions & ref.\\
\end{longtable}
}

\noindent{\scriptsize $^\dagger$A secondary summary lists $\sim$4{,}900 logical qubits; not confirmed in the primary source and not used.}

\subsection{Further figures for the chemistry, lattice and gauge cases}
Fig.~\ref{fig:femoco} shows the frontier for the smaller FeMoco model and its history. Fig.~\ref{fig:sqd} gives the sampling limits of quantum-selected configuration interaction discussed in Section~\ref{sec:devices}. Fig.~\ref{fig:dmrg} shows how DMRG cost grows with bond dimension for the P450 models: with CPU time growing as the cube of the bond dimension, a value of 12{,}000 would cost \planCoreH{} core-hours and about \planMemTB\,TB of memory, a budget available at supercomputer scale, although whether that bond dimension converges the spin gap is a separate and empty evidence field. Fig.~\ref{fig:platform} shows how one fixed algorithm depends on the hardware scenario, and Fig.~\ref{fig:lgt} the spread of gauge-theory estimates.

\begin{figure}[htbp]
\centering
\begin{tikzpicture}
\begin{groupplot}[group style={group size=2 by 1, horizontal sep=1.9cm}, height=6.0cm, ticklabel style={font=\footnotesize}, label style={font=\footnotesize}]
\nextgroupplot[width=0.56\textwidth, xmode=log, ymode=log, xmin=4e4, xmax=1.5e7, ymin=0.3, ymax=5000, xlabel={Physical qubits}, ylabel={Expected runtime (hours)}, grid=major, grid style={black!10}, clip=false]
\addplot[only marks, mark=square*, mark size=2.6pt, color=black!55] table[x=x,y=y] {figdata/femoco54_prior.dat};
\addplot[only marks, mark=*, mark size=2.8pt, color=qcF] table[x=x,y=y] {figdata/femoco54_dense.dat};
\addplot[qcF!60, dashed] table[x=x,y=y] {figdata/femoco54_dense.dat};
\node[font=\scriptsize, anchor=west] at (axis cs:6.3e6,1440) {DF};
\node[font=\scriptsize, anchor=west] at (axis cs:4.2e6,72) {THC};
\node[font=\scriptsize, anchor=west] at (axis cs:3.4e6,3) {DFTHC+SA};
\node[font=\scriptsize, anchor=south west, qcF] at (axis cs:9.5e4,720) {compact};
\node[font=\scriptsize, anchor=north east, qcF] at (axis cs:1.9e5,20) {balanced};
\node[font=\scriptsize, anchor=north east, qcF] at (axis cs:1.0e6,1.0) {efficient};
\node[font=\scriptsize, anchor=north west, qcF] at (axis cs:2.1e6,0.75) {fast};
\draw[{Latex[length=2mm]}-{Latex[length=2mm]}, black!50] (axis cs:88840,2.2) -- node[below, font=\scriptsize]{36$\times$ fewer qubits} (axis cs:3.2e6,2.2);

\nextgroupplot[width=0.40\textwidth, ybar, ymode=log, log origin=infty, ymin=0.5, ymax=2e4, xtick={1,...,7}, xticklabels={DF,THC,DFTHC,compact,balanced,efficient,fast}, xmin=0.4, xmax=7.6, xticklabel style={rotate=45, anchor=east, font=\scriptsize}, ylabel={Space-time volume ($10^6$ qubit-h)}, bar width=8pt, bar shift=0pt, nodes near coords, point meta=rawy, every node near coord/.append style={font=\tiny}]
\addplot[fill=black!20, draw=black!75, line width=0.3pt, postaction={pattern=north east lines, pattern color=black!80}] table[x=x,y=y] {figdata/femoco54vol_prior.dat};
\addplot[fill=qcF, draw=none] table[x=x,y=y] {figdata/femoco54vol_dense.dat};
\end{groupplot}
\end{tikzpicture}
\caption[FeMoco, 54 orbitals: quantum cost frontier]{Phase estimation on the 54-orbital FeMoco model at 1\,mHa, all values from one source~\cite{low2026dense}. Squares, and hatched bars in the right panel: prior estimates as tabulated there, by double factorization (DF), tensor hypercontraction (THC) and THC with spectrum amplification (DFTHC+SA). Circles and plain bars: the 2026 dense planar surface code at a 1\,$\mu$s cycle and 10\,$\mu$s reaction time, in four rows of its Table~XXVII that the source names compact (within a month), balanced (within a day), efficient (minimum volume) and fast. The source's summary table gives the first two as upper limits (under a month, under a day); the tabulated expected runtimes are plotted here. Against the best prior point the source reports \femoSpace$\times$ fewer physical qubits and \femoVol$\times$ smaller minimum space-time volume. Larger ratios against the oldest point describe progress across algorithms, not the frontier at a fixed qubit budget.}
\label{fig:femoco}
\end{figure}

\begin{figure}[htbp]
\centering
\begin{tikzpicture}
\begin{groupplot}[group style={group size=2 by 1, horizontal sep=2.0cm}, width=0.43\textwidth, height=5.4cm, ticklabel style={font=\footnotesize}, label style={font=\footnotesize}, title style={font=\footnotesize, align=center}]
\nextgroupplot[ybar, ymin=0, ymax=60000, xtick={1,2}, xticklabels={QSCI, HCI}, xmin=0.3, xmax=2.7, bar width=22pt, ylabel={Determinants}, title={[2Fe-2S] at 0.1\,Ha ($\approx$63\,kcal/mol):\\ideal QSCI vs HCI}, nodes near coords, every node near coord/.append style={font=\scriptsize}, scaled y ticks=false, yticklabel style={/pgf/number format/fixed}]
\addplot[fill=qcF, draw=none] table[x=x,y=y] {figdata/sqd_dets.dat};
\nextgroupplot[xmode=log, ymode=log, xmin=1e2, xmax=1e15, ymin=1e-6, ymax=1, xlabel={Shots}, ylabel={New determinants per shot}, title={N$_2$ CAS(10e,12o) discovery rate}, grid=major, grid style={black!10}]
\addplot[only marks, mark=*, color=qcF, mark size=2.5pt, nodes near coords={\pgfmathprintnumber[sci,precision=1]{\pgfplotspointmeta}}, point meta=rawy, every node near coord/.append style={font=\scriptsize, anchor=south west}] table[x=x,y=y] {figdata/sqd_rate.dat};
\addplot[dashed, black!40] table[x=x,y=y] {figdata/sqd_interp.dat};
\draw[qcU, line width=0.9pt] (axis cs:4e14,1e-6) -- (axis cs:4e14,1);
\node[font=\scriptsize, anchor=north east, align=right, qcUd] at (axis cs:3e14,0.5) {$4\times10^{14}$ samples for\\ $\mu$Ha precision\\ (Reinholdt et al.)};

\end{groupplot}
\end{tikzpicture}
\caption[QSCI compactness and sampling]{Quantum-selected CI against classical selected CI, from Reinholdt et al.~\cite{reinholdt2025} (values in \texttt{values.csv}). Left: determinant counts at a 0.1\,Ha error ($\approx$63\,kcal/mol) for ideal ground-state-distribution QSCI versus heat-bath CI with a chemically informed start, in individually selected determinants. This is not a chemical-accuracy comparison and not the count of the hardware SQD experiment, which uses spin-string product subspaces. Right: determinant discovery per shot for N$_2$ (two published points; the dashed line only interpolates between them). The vertical line marks the source's own estimate of $4\times10^{14}$ samples for micro-hartree precision; the discovery-rate curve does not by itself determine energy error, and micro-hartree precision is one example and not the requirement of every application.}
\label{fig:sqd}
\end{figure}

\begin{figure}[htbp]
\centering
\begin{tikzpicture}
\begin{axis}[
  xmode=log, ymode=log, width=0.72\textwidth, height=6.2cm,
  xmin=1000, xmax=30000, ymin=1e3, ymax=1e8, xtick={1000,2000,5000,10000,20000}, xticklabels={1000,2000,5000,10000,20000},
  xlabel={\footnotesize Bond dimension $M$ (or $D$)}, ylabel={\footnotesize CPU core-hours},
  ticklabel style={font=\footnotesize}, grid=major, grid style={black!10},
  legend style={font=\scriptsize, at={(0.98,0.02)}, anchor=south east, draw=black!25, fill=white}, legend cell align=left]
\addplot[only marks, mark=square*, mark size=2.6pt, color=qcF] table[x=x,y=y] {figdata/dmrg_measured_G.dat};
\addlegendentry{P450 model G, 41--43 orbitals (measured)}
\addplot[only marks, mark=o, mark size=2.8pt, color=qcF, line width=1pt] table[x=x,y=y] {figdata/dmrg_estimated_X.dat};
\addlegendentry{P450 model X, 58 orbitals (estimated in source)}
\addplot[qcF, dashed] table[x=M,y=cpuh] {figdata/dmrg_plan.dat};
\addlegendentry{$M^3$ scaling, planning scenario}
\addplot[only marks, mark=triangle*, mark size=3pt, color=black!65] table[x=x,y=y] {figdata/dmrg_femoco_sweep.dat};
\addlegendentry{FeMoco-76 UDMRG, one sweep (measured)}
\draw[black!40, dashed] (axis cs:1000,7.63e6) -- (axis cs:30000,7.63e6) node[pos=0.03, above, anchor=south west, font=\scriptsize, black!60]{allocation equivalent: all Fugaku cores for one hour (not a runtime)};

\end{axis}
\end{tikzpicture}
\caption[DMRG cost versus bond dimension]{DMRG cost for the P450 active spaces~\cite{goings2022}. Only model G was measured. Both model-X values are the source's own estimates from theoretical scaling, so they do not independently validate $M^3$ scaling, and the source gives no guarantee that $M=3000$ suffices. The dashed curve extends that scaling as a resource-planning scenario and is not evidence of a converged gap. The FeMoco-76 sweep~\cite{zhai2026} shows that budgets of this size are affordable; it does not show that a P450 spin gap converges at any particular $M$, since bond-dimension requirements are basis- and problem-dependent.}
\label{fig:dmrg}
\end{figure}

\begin{figure}[htbp]
\centering
\begin{tikzpicture}
\begin{axis}[
  ybar, ymode=log, log origin=infty, width=0.76\textwidth, height=6.0cm,
  ymin=0.1, ymax=5e10, xtick={1,2,3,4}, xticklabels={{$\mu$s, $p{=}10^{-3}$}, {$\mu$s, $p{=}10^{-4}$}, {ns, $p{=}10^{-3}$}, {ns, $p{=}10^{-4}$}}, xmin=0.4, xmax=4.6,
  ylabel={\footnotesize Runtime per execution (s)}, ticklabel style={font=\footnotesize}, bar width=14pt,
  legend style={font=\scriptsize, at={(0.5,1.03)}, anchor=south, draw=black!25, legend columns=1},
  ymajorgrids, grid style={black!10}]
\addplot[fill=qcF, draw=none] table[x=x,y=y] {figdata/platform_TFIM.dat};
\addlegendentry{2D TFIM dynamics (230 LQ, $2.4\times10^6$ T)}
\addplot[fill=black!20, draw=black!75, line width=0.3pt, postaction={pattern=north east lines, pattern color=black!80}] table[x=x,y=y] {figdata/platform_RU.dat};
\addlegendentry{Ru catalyst (2740 LQ, $5.4\times10^{11}$ T)}
\end{axis}
\end{tikzpicture}
\caption[Source schedule sensitivity]{Source schedule sensitivity: one fixed algorithm per target under four qubit scenarios from Beverland et al.~\cite{beverland2022}. Each scenario has its own factory allocation and code distance as set by the source, so this compares source schedules and does not isolate gate speed. For the catalyst the runtime spans a factor of about \ruSpread{} across the four scenarios and about \ruSpreadFixed{} between the two gate platforms at a fixed physical error rate.}
\label{fig:platform}
\end{figure}

\begin{figure}[htbp]
\centering
\begin{tikzpicture}
\begin{axis}[
  xmode=log, width=0.70\textwidth, height=4.4cm, title={\footnotesize Different approximations and error criteria: not a progress series}, title style={font=\footnotesize},
  xmin=1e7, xmax=1e52, ymin=0.4, ymax=3.6, y dir=reverse, ytick={1,...,3},
  yticklabels={{Electric-basis Trotter (Kan--Nam type)}, {Electric-basis LCU}, {Large-$N_c$/Krylov truncations}},
  yticklabel style={font=\footnotesize}, xlabel={\footnotesize T gates for time evolution only},
  xticklabel style={font=\footnotesize},
  axis x line*=bottom, axis y line*=left, xmajorgrids, grid style={black!10}, clip=false]
\addplot[only marks, mark=|, mark size=5pt, black!65, line width=1.5pt] table[x=x,y=y] {figdata/gauge_lo.dat};
\addplot[only marks, mark=|, mark size=5pt, black!65, line width=1.5pt] table[x=x,y=y] {figdata/gauge_hi.dat};
\node[font=\scriptsize, anchor=east, inner sep=5pt] at (axis cs:6.5e48,1) {$6.5\times10^{48}$ T, $\sim10^5$ LQ};
\draw[black!60, line width=4pt] (axis cs:1e11,2) -- (axis cs:1e14,2);
\node[font=\scriptsize, anchor=west, inner sep=5pt] at (axis cs:1e14,2) {$10^{11}$--$10^{14}$ T};
\draw[black!60, line width=4pt] (axis cs:1e9,3) -- (axis cs:1e12,3);
\node[font=\scriptsize, anchor=west, inner sep=5pt] at (axis cs:1e12,3) {$10^{9}$--$10^{12}$ T};

\end{axis}
\end{tikzpicture}
\caption[Lattice gauge theory estimates by formulation]{Published T-gate estimates for time evolution in a fiducial real-time SU(3) shear-viscosity calculation (pure gauge, $10^3$ spatial lattice), one row per formulation, from the compilation by Lamm~\cite{lamm2026} with the truncations of Ciavarella et al.~\cite{ciavarella2025}. The $\sim10^5$ logical-qubit width of the first row is as cited by Carena et al.~\cite{carena2022}; the other rows report no width. The rows use different approximations and error criteria and are \emph{not} a progress series at a common controlled accuracy: several aggressive truncations are not shown to support continuum extrapolation, and synthesis precision is not the error in a transport coefficient. All exclude state preparation, correlator measurement and observable extraction. Two discrete-subgroup estimates near $10^{12}$ T~\cite{lamm2026,osorio2025} are omitted because we have not read them at table level.}
\label{fig:lgt}
\end{figure}


\subsection{The query-complexity result in detail}
\label{app:query}
\textbf{Statement.} Let $f$ be a real polynomial of degree at most $d\ge1$ in $N$ variables with $0\le f(x)\le1$ for all $x\in\{0,1\}^N$, and let its variance under the uniform measure be at least $\epsilon$. Then some variable has influence at least a fixed polynomial in $\epsilon/d$, with influence defined as in~\cite{aa2014}. This is Conjecture~1.7 of Aaronson and Ambainis; a proof was posted by Turturean on 6 November 2026~\cite{turturean2026,psi2026}.

\textbf{Consequence.} By a theorem of the 2014 paper, the acceptance probability of any quantum algorithm making $T$ queries can then be approximated to within $\epsilon$, on all but a fraction $\delta$ of inputs, by a deterministic classical algorithm making at most $(4T)^8\epsilon^{-14}\delta^{-8}$ queries~\cite{turturean2026}. An earlier special case covers algorithms with a bounded number of parallel query rounds~\cite{bdst2026}, and equivalent statements hold for biased product and slice distributions~\cite{gmnr2026}.

\textbf{Verification status.} Three layers. (i) The influence inequality is formalized in Lean~4; the repository reports no unfinished steps and no added axioms in the library, at commit \texttt{2a6ceae}. We read the formal statement, which matches the conjecture as printed, and did not rebuild the proof. (ii) That the formal statement is the conjecture is a human judgment; an expert exposition exists~\cite{blancexpo2026}, and Aaronson states that he believes the proof correct~\cite{aaronsonblog2026}. (iii) The query consequence follows through the peer-reviewed 2014 theorem and is not part of the formal development. The work is a preprint, and an earlier proof of the same conjecture, announced in 2019~\cite{aaronson2019}, was later found to contain an error. We record it as a claimed proof with a machine-checked core.

\textbf{What ``polynomial'' is worth.} With $T=10^3$ and $\epsilon=\delta=0.1$ the bound is \aaBound{} classical queries. The purpose of the example is to separate an asymptotic upper bound from an operational cost: the theorem excludes an exponential separation in queries and says nothing about a practical crossover. It does not enter the break-even model.

\textbf{How it was produced.} According to the author, the mathematical statements, the first draft and the formalization were produced by AI models in a research harness, using a non-public model; the author scoped the problem, reviewed the paper and takes responsibility for it~\cite{turturean2026,psi2026}. The mathematical artifact can be checked by anyone. The reported discovery run cannot be independently reproduced from the released artifacts, and the separate contributions of harness, model and author are not isolated.
